% !TEX root =./ECE444_Practice_main.tex
%
% L17 practice set (Introduction to Phased Array Hardware) -- identification of
% the ADALM-PHASER receive chain, the fixed-IF frequency plan, and the
% consequences of hybrid (analog + digital) beamforming.

\fullwidth{\textbf{Learning Objective 3.3: } I can identify the hardware architecture
of the ADALM-PHASER and control it via SDR software.
\\[4pt]\normalfont\small Show your work. A \textbf{Documentation} line at the end
is required (write \textbf{None} if you did not collaborate).}

\begin{questions}

\question \textbf{Walking the receive chain.} A wave from an X-band source arrives at
the PHASER's patch face and ends up as numbers in the Raspberry Pi.

\begin{parts}

	\part List, in order, the blocks the signal passes through from a single patch to the
	Raspberry Pi. Name the part number or device for each one.
		\begin{solution}[1.3in]
		Patch element, then the ADL8107 LNA behind it, then the ADAR1000 analog
		beamformer, which applies phase and gain and sums its four elements. That sum
		goes to an LTC5548 mixer driven by the ADF4159 and HMC735 LO, which produces the
		2.2 GHz IF. The IF goes to one receive channel of the ADALM-Pluto (AD9361), and
		the Raspberry Pi reads the IQ samples and serves the browser interface.
		\end{solution}

	\begin{minipage}[t]{0.58\linewidth}
		\part The ADAR1000 sets element phase with a 7-bit phase shifter. Compute the phase
		LSB in degrees, and state what physically limits how finely you can steer the beam.
		\begin{solution}[1.1in]
		\eq{ \text{LSB} = \frac{360\mydeg}{2^{7}} = \frac{360\mydeg}{128} = 2.8125\mydeg }
		The commanded phase is rounded to the nearest multiple of this step, so the applied
		phase ramp -- and therefore the beam direction -- can only take discrete values.
		The limit is the phase-shifter word length, not the software.
		\end{solution}
	\end{minipage}\hfill%
	\begin{minipage}[t]{0.37\linewidth}
		\raggedleft \vspace{12pt}
		LSB = \ansbox[1.3in]{$2.8125\mydeg$}
	\end{minipage}

	\begin{minipage}[t]{0.58\linewidth}
		\part How many array elements feed one ADAR1000, and how many RF channels leave the
		pair of ADAR1000s on their way to the mixers?
		\begin{solution}[0.9in]
		Each ADAR1000 is a four-channel beamformer, so four of the eight elements feed each
		chip. Each chip sums its four inputs into one RF output, so two RF channels leave
		the pair -- one per 4-element subarray.
		\end{solution}
	\end{minipage}\hfill%
	\begin{minipage}[t]{0.37\linewidth}
		\raggedleft \vspace{12pt}
		elements/chip = \ansbox[0.9in]{$4$}
		\\[10pt]
		RF channels = \ansbox[0.9in]{$2$}
	\end{minipage}

	\begin{minipage}[t]{0.58\linewidth}
		\part The element spacing is $d = 14\ \text{mm}$. For a source at $10.525\ \ghz$,
		find the wavelength and the spacing in wavelengths.
		\begin{solution}[1.2in]
		\eq{ \lambda &= \frac{c}{f} = \frac{3\times10^{8}\ \text{m/s}}{10.525\times10^{9}\ \text{Hz}} = 0.0285\ \m = 28.5\ \text{mm} \\
		     \frac{d}{\lambda} &= \frac{14\ \text{mm}}{28.5\ \text{mm}} = 0.491 }
		Just under a half wavelength, which is the standard choice: close enough to
		$\lambda/2$ for a full scan range, but never above it, so no grating lobe enters
		visible space.
		\end{solution}
	\end{minipage}\hfill%
	\begin{minipage}[t]{0.37\linewidth}
		\raggedleft \vspace{12pt}
		$\lambda$ = \ansbox[1.1in]{$28.5\ \text{mm}$}
		\\[10pt]
		$d/\lambda$ = \ansbox[1.1in]{$0.491$}
	\end{minipage}

\end{parts}

\newpage

\question \textbf{The frequency plan.} The PHASER mixes the received X-band signal down
to a fixed $2.2\ \ghz$ IF using high-side injection, so $f_{\text{LO}} = f_{\text{RF}} +
f_{\text{IF}}$. The ADF4159 and HMC735 VCO cover $12.2$ to $13.0\ \ghz$. The Pluto is
tuned to $2.2\ \ghz$ and sampled at $3\ \text{MSPS}$.

\begin{parts}

	\begin{minipage}[t]{0.58\linewidth}
		\part \textbf{Find HB100} reports the lab source at $10.42\ \ghz$. What LO frequency
		must the software command?
		\begin{solution}[1.0in]
		\eq{ f_{\text{LO}} = f_{\text{RF}} + f_{\text{IF}} = 10.42 + 2.20 = 12.62\ \ghz }
		This is inside the $12.2$--$13.0\ \ghz$ VCO range, so the source is reachable.
		\end{solution}
	\end{minipage}\hfill%
	\begin{minipage}[t]{0.37\linewidth}
		\raggedleft \vspace{12pt}
		$f_{\text{LO}}$ = \ansbox[1.3in]{$12.62\ \ghz$}
	\end{minipage}

	\begin{minipage}[t]{0.58\linewidth}
		\part Given only the VCO limits and the fixed IF, what range of source frequencies
		can this receiver reach?
		\begin{solution}[1.0in]
		\eq{ f_{\text{RF}} = f_{\text{LO}} - f_{\text{IF}}: \quad
		     12.2 - 2.2 &= 10.0\ \ghz \\
		     13.0 - 2.2 &= 10.8\ \ghz }
		The receiver covers $10.0$ to $10.8\ \ghz$, which contains the whole $10.1$ to
		$10.7\ \ghz$ spread of HB100 units.
		\end{solution}
	\end{minipage}\hfill%
	\begin{minipage}[t]{0.37\linewidth}
		\raggedleft \vspace{12pt}
		$f_{\text{RF,min}}$ = \ansbox[1.1in]{$10.0\ \ghz$}
		\\[10pt]
		$f_{\text{RF,max}}$ = \ansbox[1.1in]{$10.8\ \ghz$}
	\end{minipage}

	\begin{minipage}[t]{0.58\linewidth}
		\part A student skips \textbf{Find HB100}, so the software keeps the nominal LO of
		$12.725\ \ghz$. The source in front of the array is actually at $10.30\ \ghz$. What
		IF does the mixer produce, and by how much does it miss $2.2\ \ghz$?
		\begin{solution}[1.2in]
		\eq{ f_{\text{IF}} &= f_{\text{LO}} - f_{\text{RF}} = 12.725 - 10.300 = 2.425\ \ghz \\
		     \Delta f &= 2.425 - 2.200 = 0.225\ \ghz = 225\ \mhz }
		The tone lands 225 MHz above where the Pluto is tuned.
		\end{solution}
	\end{minipage}\hfill%
	\begin{minipage}[t]{0.37\linewidth}
		\raggedleft \vspace{12pt}
		$f_{\text{IF}}$ = \ansbox[1.1in]{$2.425\ \ghz$}
		\\[10pt]
		$\Delta f$ = \ansbox[1.1in]{$225\ \mhz$}
	\end{minipage}

	\begin{minipage}[t]{0.58\linewidth}
		\part Find the width of the baseband window the Pluto digitizes at $3\ \text{MSPS}$,
		and say in one or two sentences what the student of part (c) sees on the FFT tab.
		\begin{solution}[1.2in]
		\eq{ \text{window} = f_s = 3\ \mhz \quad (\pm 1.5\ \mhz) }
		The window is centred on the 2.2 GHz tuning. The offset of 225 MHz is 150 times
		the half-width of that window, so the tone is far
		outside it and nothing appears but the noise floor. An empty FFT with a powered,
		correctly aimed source is a frequency error, not a hardware failure.
		\end{solution}
	\end{minipage}\hfill%
	\begin{minipage}[t]{0.37\linewidth}
		\raggedleft \vspace{12pt}
		window = \ansbox[1.3in]{$3\ \mhz$}
	\end{minipage}

\end{parts}

\newpage

\question \textbf{Hybrid beamforming.} The PHASER has eight elements, eight RF phase
shifters, and two receive channels into the SDR.

\begin{parts}

	\part Why does the board put a phase shifter behind every element at RF instead of
	digitizing all eight elements and doing the phase shift in software?
		\begin{solution}[1.2in]
		An RF phase shifter and attenuator are small, cheap, and low-power, so eight of them
		cost little. A full receive channel -- mixer, filter, ADC, and the data path to move
		the samples -- is expensive in parts, board area, and power. Digitizing all eight
		elements would multiply the receiver hardware by four for a teaching board that
		mostly needs one beam at a time.
		\end{solution}

	\begin{minipage}[t]{0.58\linewidth}
		\part After the two 4:1 analog sums, how many independent complex weights can
		software still apply, and what information has been destroyed?
		\begin{solution}[1.1in]
		Two, one per digital channel. Once four elements are summed in the ADAR1000 their
		individual amplitudes and phases no longer exist anywhere in the system; only the
		single sum survives. Software can weight and combine the two subarray outputs, but
		it cannot recover what any one element received.
		\end{solution}
	\end{minipage}\hfill%
	\begin{minipage}[t]{0.37\linewidth}
		\raggedleft \vspace{12pt}
		digital weights = \ansbox[1.0in]{$2$}
	\end{minipage}

	\part An adaptive nulling algorithm running on the digital channels can steer roughly
	one null per available degree of freedom. State one capability the hybrid split keeps
	and one it takes away, and say which of the two you would trade if the board's job were
	to null several jammers at once.
		\begin{solution}[1.4in]
		Kept: full per-element control of amplitude and phase at RF, so the array can steer,
		taper, and place a fixed null anywhere the designer computes in advance. Taken away:
		adaptive, data-driven nulling, because with two digital channels the adaptive
		processor has only one spare degree of freedom and can suppress about one interferer.
		To null several jammers at once you would trade the cheap analog architecture for
		more receive channels -- ideally one per element -- and accept the cost and power
		that come with it.
		\end{solution}

	\begin{minipage}[t]{0.58\linewidth}
		\part How many ADC channels would a fully digital 8-element version need, and name
		one practical cost besides part count.
		\begin{solution}[1.0in]
		Eight, one per element. Beyond the parts, the data rate off the board rises by a
		factor of four, the power budget grows with the added mixers and converters, and
		every channel needs its own amplitude and phase calibration before the samples can
		be combined coherently.
		\end{solution}
	\end{minipage}\hfill%
	\begin{minipage}[t]{0.37\linewidth}
		\raggedleft \vspace{12pt}
		ADC channels = \ansbox[1.0in]{$8$}
	\end{minipage}

\end{parts}

\newpage

\question \textbf{Reading the bench.} Lab Preset 1 is loaded, the HB100 is at boresight
about $1\ \m$ from the array, and the FFT tab is showing the baseband spectrum. The Pluto
is tuned to $2.2\ \ghz$ and the LO is at $12.725\ \ghz$.

\begin{parts}

	\begin{minipage}[t]{0.58\linewidth}
		\part The FFT peak sits $1\ \mhz$ above the centre of the window. Find the IF of the
		received tone and the frequency of the source.
		\begin{solution}[1.3in]
		\eq{ f_{\text{IF}} &= 2.200 + 0.001 = 2.201\ \ghz \\
		     f_{\text{RF}} &= f_{\text{LO}} - f_{\text{IF}} = 12.725 - 2.201 = 10.524\ \ghz }
		The source is 1 MHz below the nominal 10.525 GHz, which is ordinary for a
		free-running HB100.
		\end{solution}
	\end{minipage}\hfill%
	\begin{minipage}[t]{0.37\linewidth}
		\raggedleft \vspace{12pt}
		$f_{\text{IF}}$ = \ansbox[1.1in]{$2.201\ \ghz$}
		\\[10pt]
		$f_{\text{RF}}$ = \ansbox[1.1in]{$10.524\ \ghz$}
	\end{minipage}

	\part \textbf{Rx Gain} is raised from $10\db$ to $30\db$. Predict what happens to the
	peak level, to the noise floor, and to the separation between them, and explain why.
		\begin{solution}[1.3in]
		Both the peak and the noise floor rise by close to 20 dB, and the separation between
		them changes very little. The Rx Gain is applied inside the AD9361, well downstream
		of the LNAs that already fixed how much noise accompanies the signal, so it amplifies
		signal and noise together. The one exception is the low end: at $10\db$ the SDR's own
		noise and quantization floor make a measurable contribution, so the separation is
		slightly worse there than at $30\db$.
		\end{solution}

	\part A student reports no peak anywhere in the FFT. Give the first two checks you would
	make, in order, and say what each one rules out.
		\begin{solution}[1.2in]
		First, confirm the HB100 is powered and aimed at the array from about a metre --
		this rules out having no signal to receive at all, which is the most common cause.
		Second, run \textbf{Find HB100} and compare the result with the Signal Freq the GUI
		is using -- this rules out a source sitting outside the 3 MHz baseband window
		because the LO is set for the wrong frequency. Only after both pass is it worth
		suspecting the array or the cabling.
		\end{solution}

	\part On a uniform-taper beam sweep this board normally shows a noise floor about
	$23\db$ below the peak. A student's sweep shows only $8\db$. Give two plausible causes
	and say how you would tell them apart.
		\begin{solution}[1.4in]
		Either the received signal is too weak -- source too far away, misaimed, or its
		battery low -- so the peak has come down toward a fixed floor; or the array has not
		been calibrated, so the eight elements do not add coherently at broadside and the
		main lobe never reaches full height. To separate them, move the source closer and
		re-aim it: if the separation improves, it was signal level. If it does not, run
		\textbf{Calibrate} and sweep again.
		\end{solution}

\end{parts}

\vspace{1.5em}
\noindent\textbf{Documentation:} \rule[-2pt]{0.6\linewidth}{0.4pt}

\end{questions}
